% Options for packages loaded elsewhere
\PassOptionsToPackage{unicode}{hyperref}
\PassOptionsToPackage{hyphens}{url}
\documentclass[
]{article}
\usepackage{ma2003b-notes}
\hypersetup{
  pdftitle={Factor Analysis Study Guide},
  pdfauthor={Dr. Juliho Castillo},
  hidelinks,
  pdfcreator={LaTeX via pandoc}}

\title{Factor Analysis Study Guide}
\author{Dr. Juliho Castillo}
\date{}

\begin{document}
\maketitle

\textbf{Factor Analysis Study Guide}

Comprehensive Lecture Notes

Dr. Juliho Castillo

Tecnológico de Monterrey

13:54:59 2026-08-28

\begin{center}\rule{0.5\linewidth}{0.5pt}\end{center}

\section{Introduction}

This study guide provides comprehensive coverage of Principal Component
Analysis (PCA) and Factor Analysis, two fundamental multivariate
statistical techniques. The material is organized to help students
understand both the theoretical foundations and practical applications
of these methods.

\textbf{Learning Objectives:}

\begin{itemize}
\item
  Understand the mathematical foundations of PCA and Factor Analysis
\item
  Learn when to apply each method appropriately
\item
  Gain hands-on experience through real-world examples
\item
  Develop skills in interpreting and communicating results
\end{itemize}

\section{Theoretical Foundations}

\subsection{Principal Component Analysis (PCA) Theory}

\subsubsection{What is Principal Component Analysis?}

Principal Component Analysis (PCA) is a fundamental technique in
multivariate statistics and data science. At its core, PCA is a
\textbf{linear transformation} that reduces the dimensionality of data
while preserving as much variance as possible.

\textbf{Key Characteristics:}

\begin{itemize}
\item
  \textbf{Dimension Reduction}: Transforms high-dimensional data into a
  lower-dimensional space
\item
  \textbf{Variance Maximization}: Each principal component captures the
  maximum possible variance
\item
  \textbf{Orthogonality}: Principal components are uncorrelated
  (orthogonal to each other)
\item
  \textbf{Linear Transformation}: Components are linear combinations of
  original variables
\end{itemize}

\textbf{Primary Applications:}

\begin{enumerate}
\item
  \textbf{Data Visualization}: Reducing data to 2D or 3D for plotting
\item
  \textbf{Noise Reduction}: Filtering out low-variance components
\item
  \textbf{Feature Engineering}: Creating new features for machine
  learning
\item
  \textbf{Exploratory Data Analysis}: Understanding data structure and
  patterns
\end{enumerate}

\subsubsection{Mathematical Foundation}

Let \(\mathbf{X}\) be an \(n \times p\) data matrix where:

\begin{itemize}
\item
  \(n\) = number of observations (rows)
\item
  \(p\) = number of variables (columns)
\item
  Each row represents one observation
\item
  Each column represents one variable
\end{itemize}

\textbf{Step 1: Data Centering}

First, we center the data by subtracting the mean from each variable:
\[\mathbf{X}_{\text{centered }} = \mathbf{X} - \mathbf{1}_{n}\mathbf{\mu}^{\top}\]

where
\(\mathbf{\mu} = \left( \mu_{1},\mu_{2},\ldots,\mu_{p} \right)^{\top}\)
is the vector of variable means.

\textbf{Step 2: Covariance Matrix}

The sample covariance matrix is:
\[\mathbf{S} = \frac{1}{n - 1}\mathbf{X}_{\text{centered}}^{\top}\mathbf{X}_{\text{centered }}\]

For standardized analysis (correlation matrix), we first standardize
each variable:
\[\mathbf{R} = \frac{1}{n - 1}\mathbf{Z}^{\top}\mathbf{Z}\]

where \(\mathbf{Z}\) contains standardized variables (mean 0, variance
1).

\textbf{Step 3: Eigendecomposition}

The heart of PCA is the eigendecomposition of the covariance (or
correlation) matrix:
\[\mathbf{S}\mathbf{v}_{j} = \lambda_{j}\mathbf{v}_{j}\]

where:

\begin{itemize}
\item
  \(\lambda_{1} \geq \lambda_{2} \geq \ldots \geq \lambda_{p} \geq 0\)
  are eigenvalues
\item
  \(\mathbf{v}_{1},\mathbf{v}_{2},\ldots,\mathbf{v}_{p}\) are
  corresponding eigenvectors
\item
  Eigenvectors are orthonormal:
  \(\mathbf{v}_{i}^{\top}\mathbf{v}_{j} = \delta_{ij}\)
\end{itemize}

\textbf{Step 4: Principal Components}

The \(j\)-th principal component for observation \(i\) is:
\[\text{ PC}_{ij} = \mathbf{v}_{j}^{\top}\left( \mathbf{x}_{i} - \mathbf{\mu} \right)\]

In matrix form, the principal component scores are:
\[\mathbf{Y} = \mathbf{X}_{\text{centered }}\mathbf{V}\]

where
\(\mathbf{V} = \left\lbrack \mathbf{v}_{1},\mathbf{v}_{2},\ldots,\mathbf{v}_{p} \right\rbrack\)
is the matrix of eigenvectors.

\subsubsection{Variance Explained}

Each principal component explains a specific amount of variance:

\textbf{Individual Variance:}
\[\text{ Var}\left( \text{PC}_{j} \right) = \lambda_{j}\]

\textbf{Proportion of Variance Explained:}
\[\text{ Proportion}_{j} = \frac{\lambda_{j}}{\sum_{k = 1}^{p}\lambda_{k}}\]

\textbf{Cumulative Variance Explained:}
\[\text{ Cumulative}_{j} = \frac{\sum_{k = 1}^{j}\lambda_{k}}{\sum_{k = 1}^{p}\lambda_{k}}\]

\textbf{Important Property:} The total variance is preserved:
\(\sum_{j = 1}^{p}\lambda_{j} = \sum_{j = 1}^{p}\text{ Var}\left( X_{j} \right)\)

\subsubsection{Practical Implementation Steps}

\textbf{Step 1: Data Preparation}

\begin{verbatim}
1. Check for missing values and outliers
2. Decide on standardization:
   - Use correlation matrix if variables have different scales
   - Use covariance matrix if variables are on similar scales
3. Center (and possibly standardize) the data
\end{verbatim}

\textbf{Step 2: Compute PCA}

\begin{verbatim}
1. Calculate covariance or correlation matrix
2. Perform eigendecomposition
3. Sort eigenvalues and eigenvectors in descending order
4. Compute principal component scores
\end{verbatim}

\textbf{Step 3: Determine Number of Components}

\begin{verbatim}
1. Kaiser criterion: retain components with eigenvalue > 1
2. Scree plot: look for "elbow" where slope levels off
3. Cumulative variance: retain components explaining 70-90% of variance
4. Parallel analysis: compare to random data eigenvalues
\end{verbatim}

\textbf{Step 4: Interpretation}

\begin{verbatim}
1. Examine component loadings (eigenvectors)
2. Identify which variables contribute most to each component
3. Name components based on variable patterns
4. Validate interpretation with domain knowledge
\end{verbatim}

\subsubsection{Component Loadings and Interpretation}

The eigenvectors \(\mathbf{v}_{j}\) are the \textbf{component
coefficients} (or weights) used to build each component score from the
standardized variables. \textbf{Loadings} are a distinct, rescaled
quantity: \(\text{loading}_{ij} = v_{ij}\sqrt{\lambda_{j}}\), the
correlation between variable \(i\) and component \(j\) for standardized
data (this matches the definition used in L03). Loadings, not raw
eigenvector coefficients, are what should be interpreted and compared
across components, since they are on a correlation-like scale
(\(\lbrack - 1,1\rbrack\)) that eigenvector coefficients alone do not
have.

\textbf{Loading Interpretation:}

\begin{itemize}
\item
  Large absolute loading: variable strongly influences the component
\item
  Positive loading: variable and component move in same direction
\item
  Negative loading: variable and component move in opposite directions
\item
  Small loading: variable has little influence on the component
\end{itemize}

\textbf{Rotation Consideration:} Unlike Factor Analysis, PCA typically
does not use rotation because:

\begin{itemize}
\item
  Rotation would redistribute variance among components
\item
  The goal is maximum variance explanation, not interpretability
\item
  Each component is defined to be orthogonal and capture maximum
  residual variance
\end{itemize}

\subsubsection{Key Assumptions and Limitations}

\textbf{Assumptions (for descriptive PCA):}

\begin{enumerate}
\item
  \textbf{Linearity}: Relationships between variables are linear
\item
  \textbf{Continuous or ordinal-scale variables}: PCA is defined for any
  numeric data; it does not require normality or \(n > p\)
\end{enumerate}

\emph{Note: Descriptive PCA (finding directions of maximum variance) is
well-defined even when \(n < p\), when the correlation matrix is
rank-deficient, or when data are non-normal -- it is simply an
eigendecomposition. \(n > p\), normality, and no perfect
multicollinearity matter for \textbf{specific} goals built on top of PCA
(e.g. certain significance tests for eigenvalues, or a unique full-rank
solution), not for PCA's existence.}

\textbf{Limitations:}

\begin{enumerate}
\item
  \textbf{Interpretability}: Components may not have clear substantive
  meaning
\item
  \textbf{All Variance}: Includes both signal and noise variance
\item
  \textbf{Linear Only}: Cannot capture nonlinear relationships
\item
  \textbf{Outlier Sensitivity}: Sensitive to extreme values
\item
  \textbf{Scale Dependency}: Results depend on whether data is
  standardized
\end{enumerate}

---

\textbf{Study Questions for Section 2.1:}

\begin{enumerate}
\item
  What is the fundamental goal of PCA?
\item
  Why do we center the data before computing PCA?
\item
  What does it mean for principal components to be orthogonal?
\item
  How do you interpret a component loading of -0.8 for a variable?
\item
  When would you use the correlation matrix vs. covariance matrix?
\end{enumerate}

\subsection{Factor Analysis Theory}

\subsubsection{What is Factor Analysis?}

Factor Analysis (FA) is a statistical method designed to model the
relationships among observed variables through a smaller number of
unobserved variables called \textbf{latent factors} or \textbf{common
factors}. Unlike PCA, which focuses on variance maximization, Factor
Analysis explicitly models the underlying structure that generates the
observed correlations.

\textbf{Fundamental Concept:} Factor Analysis assumes that observed
variables are \textbf{manifestations} of underlying latent constructs.
These latent factors cannot be directly measured but influence multiple
observable variables.

\textbf{Key Distinctions from PCA:}

\begin{itemize}
\item
  \textbf{Purpose}: Models latent constructs rather than reducing
  dimensions
\item
  \textbf{Variance}: Separates common variance from unique variance
\item
  \textbf{Error Modeling}: Models unique variance, which combines
  measurement error and variable-specific true variance (it does not
  isolate error on its own)
\item
  \textbf{Exploratory vs. Confirmatory}: What this chapter covers is
  \textbf{exploratory} factor analysis (EFA) -- it discovers a plausible
  factor structure from the data without assuming a specific number or
  pattern of loadings in advance. Testing a specific, pre-specified
  theoretical structure against the data (and evaluating its fit) is
  \textbf{confirmatory} factor analysis (CFA), a distinct method covered
  later in this chapter.
\end{itemize}

\textbf{Primary Applications:}

\begin{enumerate}
\item
  \textbf{Psychological Testing}: Measuring intelligence, personality
  traits, attitudes
\item
  \textbf{Market Research}: Understanding consumer preferences and
  behavior patterns
\item
  \textbf{Educational Assessment}: Identifying academic ability factors
\item
  \textbf{Medical Research}: Discovering disease symptom clusters
\item
  \textbf{Social Sciences}: Measuring abstract concepts like
  socioeconomic status
\end{enumerate}

\subsubsection{The Factor Analysis Model}

The core Factor Analysis model expresses each observed variable as a
linear combination of common factors plus a unique component:

\[X_{i} = \lambda_{i,1}F_{1} + \lambda_{i,2}F_{2} + \ldots + \lambda_{i,k}F_{k} + U_{i}\]

where:

\begin{itemize}
\item
  \(X_{i}\) = \(i\)-th observed variable (standardized)
\item
  \(F_{j}\) = \(j\)-th common factor (\(j = 1,2,\ldots,k\))
\item
  \(\lambda_{i,j}\) = factor loading of variable \(i\) on factor \(j\)
\item
  \(U_{i}\) = unique factor for variable \(i\)
\item
  \(k\) = number of common factors (typically \(k \ll p\))
\end{itemize}

\textbf{In Matrix Form:}
\[\mathbf{X} = \mathbf{\Lambda}\mathbf{F} + \mathbf{U}\]

where:

\begin{itemize}
\item
  \(\mathbf{X}\) = \(p \times 1\) vector of observed variables
\item
  \(\mathbf{\Lambda}\) = \(p \times k\) matrix of factor loadings
\item
  \(\mathbf{F}\) = \(k \times 1\) vector of common factors
\item
  \(\mathbf{U}\) = \(p \times 1\) vector of unique factors
\end{itemize}

\subsubsection{Model Assumptions}

\textbf{Assumptions about Factors:}

\begin{enumerate}
\item
  \(E\left\lbrack \mathbf{F} \right\rbrack = \mathbf{0}\) (factors have
  zero mean)
\item
  \(\text{Var}\left( \mathbf{F} \right) = \mathbf{I}\) (factors are
  standardized and uncorrelated)
\item
  \(E\left\lbrack \mathbf{U} \right\rbrack = \mathbf{0}\) (unique
  factors have zero mean)
\item
  \(\text{Cov}\left( \mathbf{F},\mathbf{U} \right) = \mathbf{0}\)
  (common and unique factors are uncorrelated)
\item
  \(\text{Cov}\left( \mathbf{U} \right) = \mathbf{\Psi}\) (unique
  factors may be correlated in some models)
\end{enumerate}

\textbf{Classical Factor Model Assumptions:} In the classical model,
unique factors are also assumed uncorrelated:
\[\text{ Cov}\left( U_{i},U_{j} \right) = 0\text{ for }i \neq j\]

This means
\(\mathbf{\Psi} = \text{ diag}\left( \psi_{1},\psi_{2},\ldots,\psi_{p} \right)\).

\subsubsection{Variance Decomposition}

Factor Analysis provides a fundamental decomposition of each variable's
variance:

\textbf{Total Variance = Common Variance + Unique Variance}

For variable \(i\):
\[\text{ Var}\left( X_{i} \right) = \sum_{j = 1}^{k}\lambda_{i,j}^{2} + \psi_{i} = h_{i}^{2} + \psi_{i}\]

where:

\begin{itemize}
\item
  \(h_{i}^{2} = \sum_{j = 1}^{k}\lambda_{i,j}^{2}\) =
  \textbf{communality} (variance explained by common factors)
\item
  \(\psi_{i}\) = \textbf{unique variance} (specific variance + error
  variance)
\end{itemize}

\textbf{Key Concepts:}

\begin{itemize}
\item
  \textbf{Communality} (\(h_{i}^{2}\)): Proportion of variable's
  variance explained by common factors
\item
  \textbf{Uniqueness} (\(\psi_{i}\)): Proportion of variance not
  explained by common factors
\item
  \textbf{Specific Variance}: Systematic variance unique to one variable
\item
  \textbf{Error Variance}: Random measurement error
\end{itemize}

\subsubsection{The Correlation Structure}

Under the factor model, the correlation matrix has a specific structure:

\[\mathbf{R} = \mathbf{\Lambda}\mathbf{\Lambda}^{\top} + \mathbf{\Psi}\]

This fundamental equation shows that:

\begin{itemize}
\item
  Correlations between variables arise from shared common factors
\item
  The diagonal elements equal 1.0 (correlation of variable with itself)
\item
  Off-diagonal elements depend on factor loadings
\end{itemize}

\textbf{Reduced Correlation Matrix:} The \textbf{reduced correlation
matrix} replaces the diagonal 1's with communality estimates:
\[\mathbf{R}^{\ast} = \mathbf{\Lambda}\mathbf{\Lambda}^{\top}\]

This matrix contains only the variance attributable to common factors.

\subsubsection{Factor Extraction Methods}

Several methods exist for estimating factor loadings:

\textbf{1. Principal Axis Factoring (PAF)}

\begin{itemize}
\item
  Most common method
\item
  Iteratively estimates communalities
\item
  Uses reduced correlation matrix
\item
  Procedure:

  \begin{enumerate}
  \item
    Start with initial communality estimates
  \item
    Perform eigendecomposition of reduced correlation matrix
  \item
    Extract factors and update communality estimates
  \item
    Repeat until convergence
  \end{enumerate}
\end{itemize}

\textbf{2. Maximum Likelihood (ML)}

\begin{itemize}
\item
  Assumes multivariate normal distribution
\item
  Provides statistical tests for number of factors
\item
  Allows goodness-of-fit testing
\item
  Computationally more intensive
\end{itemize}

\textbf{3. Principal Components Method}

\begin{itemize}
\item
  Uses principal components as initial solution
\item
  Simpler but less theoretically justified
\item
  May overextract factors
\end{itemize}

\textbf{4. Minimum Residual (MINRES)}

\begin{itemize}
\item
  Minimizes squared residuals
\item
  Good compromise between PAF and ML
\item
  Robust to distributional assumptions
\end{itemize}

\subsubsection{Determining the Number of Factors}

\textbf{1. Kaiser Criterion}

\begin{itemize}
\item
  Retain factors with eigenvalues \textgreater{} 1.0
\item
  Applied to the eigenvalues of the \textbf{original} correlation matrix
  (diagonal of 1's), not the reduced/common-factor correlation matrix
  used for extraction -- this matches standard software implementations
  (including \texttt{factor\_analyzer.get\_eigenvalues()}, which returns
  both)
\item
  May overextract in some cases
\end{itemize}

\textbf{2. Scree Test}

\begin{itemize}
\item
  Plot eigenvalues and look for ``elbow''
\item
  Retain factors before the sharp drop
\item
  Subjective but often effective
\end{itemize}

\textbf{3. Parallel Analysis}

\begin{itemize}
\item
  Compare eigenvalues to those from random data
\item
  More accurate than Kaiser criterion
\item
  Accounts for sampling variability
\end{itemize}

\textbf{4. Theoretical Considerations}

\begin{itemize}
\item
  Based on substantive theory
\item
  Number of expected constructs
\item
  Interpretability of factors
\end{itemize}

\textbf{5. Goodness-of-Fit Tests (ML only)}

\begin{itemize}
\item
  Chi-square test for model fit
\item
  Various fit indices (CFI, TLI, RMSEA)
\item
  Formal statistical criteria
\end{itemize}

\subsubsection{Factor Rotation}

Raw factor solutions are often difficult to interpret. \textbf{Factor
rotation} seeks a more interpretable solution while maintaining the same
overall fit.

\textbf{Goals of Rotation:}

\begin{itemize}
\item
  Achieve ``simple structure''
\item
  Each variable loads highly on few factors
\item
  Each factor is defined by several variables
\item
  Minimize cross-loadings
\end{itemize}

\textbf{Types of Rotation:}

\textbf{Orthogonal Rotation} (factors remain uncorrelated):

\begin{itemize}
\item
  \textbf{Varimax}: Maximizes variance of squared loadings within
  factors
\item
  \textbf{Quartimax}: Minimizes number of factors needed to explain each
  variable
\item
  \textbf{Equamax}: Compromise between Varimax and Quartimax
\end{itemize}

\textbf{Oblique Rotation} (factors allowed to correlate):

\begin{itemize}
\item
  \textbf{Promax}: Oblique version of Varimax
\item
  \textbf{Direct Oblimin}: Controls degree of correlation between
  factors
\item
  \textbf{Harris-Kaiser}: Orthogonal followed by oblique transformation
\end{itemize}

\textbf{Choosing Rotation Method:}

\begin{itemize}
\item
  Use orthogonal if factors should be uncorrelated theoretically
\item
  Use oblique if factors are expected to correlate
\item
  Varimax is most popular for orthogonal rotation
\item
  Promax is most popular for oblique rotation
\end{itemize}

\subsubsection{Interpretation Guidelines}

\textbf{Factor Loading Interpretation:}

\begin{itemize}
\item
  \(|\lambda_{i,j}| \geq 0.70\): Excellent indicator of factor
\item
  \(|\lambda_{i,j}| \geq 0.60\): Good indicator of factor
\item
  \(|\lambda_{i,j}| \geq 0.50\): Fair indicator of factor
\item
  \(|\lambda_{i,j}| \geq 0.40\): Poor but possibly significant
\item
  \(|\lambda_{i,j}| < 0.40\): Generally not interpretable
\end{itemize}

\textbf{Cross-Loadings:}

\begin{itemize}
\item
  Variables should load highly on one factor
\item
  Cross-loadings \textgreater{} 0.40 may indicate:

  \begin{itemize}
  \item
    Variable measures multiple constructs
  \item
    Need for additional factors
  \item
    Poor item design
  \end{itemize}
\end{itemize}

\textbf{Factor Naming:}

\begin{itemize}
\item
  Examine variables with highest loadings
\item
  Consider theoretical meaning
\item
  Use domain knowledge
\item
  Avoid over-interpretation
\end{itemize}

\subsubsection{Model Evaluation}

\textbf{Communality Assessment:}

\begin{itemize}
\item
  High communalities (\(h^{2} > 0.70\)): Variables well-explained by
  factors
\item
  Low communalities (\(h^{2} < 0.40\)): Variables poorly explained
\item
  Very low communalities may indicate:

  \begin{itemize}
  \item
    Variable doesn't belong in analysis
  \item
    Additional factors needed
  \item
    Measurement problems
  \end{itemize}
\end{itemize}

\textbf{Overall Model Fit:}

\begin{itemize}
\item
  Proportion of variance explained by common factors
\item
  Residual correlations should be small
\item
  Theoretical sensibility of factor structure
\end{itemize}

\textbf{Replication and Validation:}

\begin{itemize}
\item
  Cross-validation with independent samples
\item
  Confirmatory factor analysis
\item
  External validity checks
\end{itemize}

---

\textbf{Study Questions for Section 2.2:}

\begin{enumerate}
\item
  What is the key difference between Factor Analysis and PCA in terms of
  their goals?
\item
  Explain the concept of communality and how it differs from total
  variance.
\item
  Why is factor rotation necessary, and when would you choose oblique
  over orthogonal rotation?
\item
  How do you interpret a factor loading of 0.75 for a variable on a
  factor?
\item
  What does it mean when a variable has low communality in a factor
  analysis?
\end{enumerate}

\subsection{Practical Implementation of Factor Analysis}

\subsubsection{Step-by-Step Factor Analysis Process}

Factor Analysis involves a systematic sequence of steps, each requiring
careful consideration and validation. This section provides a
comprehensive guide to conducting Factor Analysis in practice.

\subsubsection{Step 1: Data Preparation and Assessment}

\textbf{Data Requirements:}

\begin{itemize}
\item
  \textbf{Sample Size}: Minimum 5-10 observations per variable
\item
  \textbf{Recommended}: At least 200 observations for stable results
\item
  \textbf{Variables}: Continuous or ordinal variables
\item
  \textbf{Missing Data}: Handle before analysis (listwise deletion,
  imputation)
\end{itemize}

\textbf{Data Quality Checks:}

\begin{enumerate}
\item
  \textbf{Outliers}: Identify and handle extreme values
\item
  \textbf{Normality}: Check distribution of variables (especially for ML
  estimation)
\item
  \textbf{Linearity}: Examine scatterplots between variables
\item
  \textbf{Multicollinearity}: Check for extremely high correlations (r
  \textgreater{} 0.90)
\end{enumerate}

\textbf{Variable Selection:}

\begin{itemize}
\item
  Include variables theoretically related to expected factors
\item
  Remove variables with very low correlations with others
\item
  Consider reverse-coding negatively worded items
\item
  Ensure adequate representation of each expected factor
\end{itemize}

\subsubsection{Step 2: Testing Factorability}

Before conducting Factor Analysis, verify that the data is suitable for
factoring.

\textbf{Correlation Matrix Inspection:}

\begin{itemize}
\item
  Most correlations should be ≥ 0.30
\item
  Look for patterns suggesting underlying factors
\item
  Identify variables that don't correlate with others
\end{itemize}

\textbf{Kaiser-Meyer-Olkin (KMO) Test:}
\[\text{ KMO } = \frac{\sum\sum r_{ij}^{2}}{\sum\sum r_{ij}^{2} + \sum\sum a_{ij}^{2}}\]

where \(a_{ij}\) are partial correlations.

\textbf{KMO Interpretation:}

\begin{itemize}
\item
  KMO ≥ 0.90: Excellent factorability
\item
  KMO ≥ 0.80: Good factorability
\end{itemize}

\begin{itemize}
\item
  KMO ≥ 0.70: Adequate factorability
\item
  KMO ≥ 0.60: Marginal factorability
\item
  KMO \textless{} 0.60: Poor factorability (consider removing variables)
\end{itemize}

\textbf{Bartlett's Test of Sphericity:} Tests the null hypothesis:
\(H_{0}:\mathbf{R} = \mathbf{I}\) (correlation matrix is identity)

A significant result (p \textless{} 0.05) indicates that Factor Analysis
is appropriate.

\textbf{Anti-image Correlation Matrix:}

\begin{itemize}
\item
  Examine diagonal elements (MSA values)
\item
  Variables with MSA \textless{} 0.50 should be considered for removal
\item
  Off-diagonal elements should be small
\end{itemize}

\subsubsection{Step 3: Factor Extraction}

\textbf{Choosing Extraction Method:}

\textbf{Principal Axis Factoring (PAF) - Most Common:}

\begin{verbatim}
1. Start with communality estimates in diagonal
2. Compute eigenvalues and eigenvectors of reduced R
3. Extract factors based on eigenvalues > 0
4. Compute new communality estimates
5. Repeat until convergence
\end{verbatim}

\textbf{Maximum Likelihood (ML) - When Normality Assumed:}

\begin{verbatim}
1. Assume multivariate normal distribution
2. Iteratively estimate parameters
3. Provides goodness-of-fit tests
4. Allows statistical inference
\end{verbatim}

\textbf{Implementation Considerations:}

\begin{itemize}
\item
  PAF is more robust to distributional violations
\item
  ML provides formal statistical tests
\item
  Both methods usually yield similar results
\item
  Start with PAF for exploratory analysis
\end{itemize}

\subsubsection{Step 4: Determining Number of Factors}

\textbf{Multiple Criteria Approach:}

\textbf{1. Eigenvalue Criteria:}

\begin{itemize}
\item
  Kaiser Rule: Retain factors with eigenvalues \textgreater{} 1.0
\item
  Modified Kaiser Rule: Use for correlation matrices only
\item
  Plot eigenvalues (scree plot)
\end{itemize}

\textbf{2. Parallel Analysis:}

\begin{verbatim}
For each number of factors k:
1. Generate random correlation matrices
2. Compute mean eigenvalues from random data
3. Compare actual eigenvalues to random ones
4. Retain factors where actual > random
\end{verbatim}

\textbf{3. Percentage of Variance:}

\begin{itemize}
\item
  Retain factors explaining meaningful variance
\item
  Social sciences: 50-60\% total variance
\item
  Natural sciences: 80\%+ total variance
\end{itemize}

\textbf{4. Interpretability:}

\begin{itemize}
\item
  Can factors be meaningfully named?
\item
  Do factors make theoretical sense?
\item
  Are factors sufficiently different?
\end{itemize}

\textbf{5. Statistical Tests (ML only):}

\begin{itemize}
\item
  Chi-square test for number of factors
\item
  Likelihood ratio tests
\item
  Information criteria (AIC, BIC)
\end{itemize}

\textbf{Decision Strategy:}

\begin{enumerate}
\item
  Apply multiple criteria
\item
  Consider theoretical expectations
\item
  Examine interpretability of solutions
\item
  Choose most parsimonious interpretable solution
\end{enumerate}

\subsubsection{Step 5: Factor Rotation}

\textbf{Rotation Goals:}

\begin{itemize}
\item
  Achieve simple structure
\item
  Maximize interpretability
\item
  Minimize cross-loadings
\end{itemize}

\textbf{Rotation Decision Tree:}

\begin{verbatim}
Are factors expected to be correlated?
├─ Yes → Use Oblique Rotation
│   ├─ Promax (most common)
│   ├─ Direct Oblimin
│   └─ Harris-Kaiser
└─ No → Use Orthogonal Rotation
    ├─ Varimax (most common)
    ├─ Quartimax
    └─ Equamax
\end{verbatim}

\textbf{Rotation Procedure:}

\begin{enumerate}
\item
  Start with unrotated solution
\item
  Apply chosen rotation method
\item
  Examine factor loading matrix
\item
  Check for simple structure
\item
  Consider alternative rotations if needed
\end{enumerate}

\textbf{Evaluating Rotation Quality:}

\begin{itemize}
\item
  High loadings (\textgreater0.50) on target factors
\item
  Low cross-loadings (\textless0.40)
\item
  Each factor defined by multiple variables
\item
  Factors are interpretable
\end{itemize}

\subsubsection{Step 6: Interpretation and Naming}

\textbf{Factor Interpretation Process:}

\textbf{1. Examine Factor Loadings:}

\begin{itemize}
\item
  Focus on loadings ≥ \textbar0.50\textbar{}
\item
  Consider both positive and negative loadings
\item
  Look for patterns across variables
\end{itemize}

\textbf{2. Identify Marker Variables:}

\begin{itemize}
\item
  Variables with highest loadings on each factor
\item
  Variables that load on only one factor
\item
  Use these to understand factor meaning
\end{itemize}

\textbf{3. Factor Naming Strategy:}

\begin{itemize}
\item
  Examine content of highest-loading variables
\item
  Consider theoretical framework
\item
  Use descriptive, meaningful names
\item
  Avoid over-interpretation
\end{itemize}

\textbf{Example Interpretation Process:}

\begin{verbatim}
Factor 1 Loadings:
- Math Achievement: 0.85
- Science Achievement: 0.78
- Reading Achievement: 0.72
- Vocabulary: 0.65

Interpretation: "Academic Achievement" factor
\end{verbatim}

\textbf{4. Check for Problematic Patterns:}

\begin{itemize}
\item
  Factors with few indicator variables (\textless3)
\item
  Variables that don't load substantially on any factor
\item
  Factors that are difficult to interpret
\item
  Unexpected cross-loadings
\end{itemize}

\subsubsection{Step 7: Model Evaluation and Validation}

\textbf{Adequacy of Factor Solution:}

\textbf{1. Communality Assessment:}

\begin{itemize}
\item
  Examine communalities for each variable
\item
  Low communalities (\textless0.40) indicate poor fit
\item
  Consider removing poorly explained variables
\end{itemize}

\textbf{2. Residual Analysis:}

\begin{itemize}
\item
  Examine residual correlations
\item
  Large residuals suggest additional factors needed
\item
  Systematic patterns in residuals indicate model problems
\end{itemize}

\textbf{3. Percentage of Variance Explained:}

\begin{itemize}
\item
  Total variance explained by common factors
\item
  Variance explained by each factor
\item
  Compare to established benchmarks
\end{itemize}

\textbf{4. Theoretical Consistency:}

\begin{itemize}
\item
  Do factors match theoretical expectations?
\item
  Are factor intercorrelations reasonable?
\item
  Can results be replicated?
\end{itemize}

\textbf{Validation Strategies:}

\textbf{1. Cross-Validation:}

\begin{itemize}
\item
  Split sample randomly
\item
  Conduct FA on each half
\item
  Compare factor structures
\item
  Use coefficient of congruence
\end{itemize}

\textbf{2. Confirmatory Factor Analysis:}

\begin{itemize}
\item
  Test specific factor structure
\item
  Use structural equation modeling
\item
  Assess model fit indices
\item
  Compare alternative models
\end{itemize}

\textbf{3. External Validation:}

\begin{itemize}
\item
  Correlate factors with external criteria
\item
  Examine predictive validity
\item
  Test known-groups validity
\item
  Assess convergent/discriminant validity
\end{itemize}

\subsubsection{Step 8: Reporting Results}

\textbf{Essential Elements to Report:}

\textbf{1. Sample and Variables:}

\begin{itemize}
\item
  Sample size and characteristics
\item
  Number and type of variables
\item
  Missing data handling
\end{itemize}

\textbf{2. Factorability Assessment:}

\begin{itemize}
\item
  KMO measure
\item
  Bartlett's test results
\item
  Any variables removed
\end{itemize}

\textbf{3. Extraction Details:}

\begin{itemize}
\item
  Method used (PAF, ML, etc.)
\item
  Criteria for number of factors
\item
  Total variance explained
\end{itemize}

\textbf{4. Rotation Information:}

\begin{itemize}
\item
  Type of rotation used
\item
  Rationale for choice
\end{itemize}

\textbf{5. Factor Structure:}

\begin{itemize}
\item
  Factor loading matrix
\item
  Factor correlations (if oblique)
\item
  Communalities
\end{itemize}

\textbf{6. Interpretation:}

\begin{itemize}
\item
  Factor names and descriptions
\item
  Theoretical implications
\item
  Limitations and assumptions
\end{itemize}

\textbf{Common Reporting Mistakes to Avoid:}

\begin{itemize}
\item
  Reporting only rotated or only unrotated loadings
\item
  Not reporting communalities
\item
  Insufficient detail about methodology
\item
  Over-interpreting small loadings
\item
  Not discussing limitations
\end{itemize}

---

\textbf{Study Questions for Section 2.3:}

\begin{enumerate}
\item
  What is the minimum recommended sample size for Factor Analysis and
  why?
\item
  How do you interpret a KMO value of 0.65?
\item
  What is the difference between Bartlett's test and the KMO test?
\item
  When would you choose Maximum Likelihood over Principal Axis
  Factoring?
\item
  What makes a factor solution ``interpretable''?
\item
  How do you validate a factor analysis solution?
\end{enumerate}

\subsection{Applied Examples and Case Studies}

This section presents detailed examples of Factor Analysis applications
across different domains, illustrating the complete analytical process
from data preparation to interpretation.

\subsubsection{Case Study 1: Psychological Assessment - Big Five
Personality}

\textbf{Background:} A researcher wants to validate the Big Five
personality model using a 25-item personality questionnaire administered
to 500 adults. Each item is rated on a 5-point Likert scale.

\textbf{Research Questions:}

\begin{enumerate}
\item
  Do the items cluster into five distinct personality factors?
\item
  How well do the factors correspond to theoretical expectations?
\item
  What is the reliability of each factor?
\end{enumerate}

\textbf{Step-by-Step Analysis:}

\textbf{Data Preparation:}

\begin{itemize}
\item
  Sample: N = 500 adults
\item
  Variables: 25 personality items (5 items per expected factor)
\item
  Scale: 1 = Strongly Disagree to 5 = Strongly Agree
\item
  Missing data: 2.3\% handled via listwise deletion
\end{itemize}

\textbf{Factorability Assessment:}

\begin{itemize}
\item
  KMO = 0.87 (Good factorability)
\item
  Bartlett's test: χ² = 4,832.5, df = 300, p \textless{} 0.001
\item
  All individual MSA values \textgreater{} 0.75
\end{itemize}

\textbf{Factor Extraction:}

\begin{itemize}
\item
  Method: Principal Axis Factoring
\item
  Initial eigenvalues: 6.84, 3.21, 2.67, 1.89, 1.43, 0.89, 0.76\ldots{}
\item
  Parallel analysis suggests 5 factors
\item
  Five factors explain 64.2\% of total variance
\end{itemize}

\textbf{Factor Rotation:}

\begin{itemize}
\item
  Method: Promax (oblique) - personality factors expected to correlate
\item
  Factor correlations range from 0.12 to 0.38
\end{itemize}

\textbf{Factor Structure and Interpretation:}

\textbf{Factor 1: Extraversion (15.2\% variance)}

\begin{verbatim}
Item                                    Loading
"I am the life of the party"              0.78
"I talk to a lot of different people"     0.74
"I feel comfortable around people"        0.71
"I start conversations"                   0.68
"I make friends easily"                   0.66
\end{verbatim}

\textbf{Factor 2: Conscientiousness (12.8\% variance)}

\begin{verbatim}
Item                                    Loading
"I am always prepared"                    0.81
"I pay attention to details"              0.76
"I get chores done right away"            0.73
"I like order"                           0.69
"I follow a schedule"                    0.65
\end{verbatim}

\textbf{Factor 3: Neuroticism (11.7\% variance)}

\begin{verbatim}
Item                                    Loading
"I get stressed out easily"               0.79
"I worry about things"                    0.75
"I get irritated easily"                  0.71
"I often feel blue"                      0.68
"I panic easily"                         0.63
\end{verbatim}

\textbf{Factor 4: Agreeableness (12.3\% variance)}

\begin{verbatim}
Item                                    Loading
"I sympathize with others' feelings"      0.77
"I have a soft heart"                     0.74
"I take time out for others"              0.70
"I feel others' emotions"                 0.67
"I make people feel at ease"              0.64
\end{verbatim}

\textbf{Factor 5: Openness (12.2\% variance)}

\begin{verbatim}
Item                                    Loading
"I have a vivid imagination"              0.76
"I have excellent ideas"                  0.73
"I am quick to understand things"         0.69
"I use difficult words"                   0.66
"I spend time reflecting on things"       0.62
\end{verbatim}

\textbf{Results Summary:}

\begin{itemize}
\item
  All communalities \textgreater{} 0.45
\item
  No problematic cross-loadings (\textgreater{} 0.40)
\item
  Factor structure matches theoretical expectations
\item
  Coefficient alpha reliability: 0.82-0.89 for all factors
\end{itemize}

\textbf{Interpretation:} The analysis successfully replicated the Big
Five personality structure with strong factor loadings and good
reliability. The moderate factor intercorrelations (0.12-0.38) support
the use of oblique rotation and align with personality theory.

\subsubsection{Case Study 2: Market Research - Consumer Technology
Preferences}

\textbf{Background:} A technology company surveyed 800 consumers about
their preferences for smartphone features using 20 attributes rated on
importance (1-7 scale).

\textbf{Research Objective:} Identify underlying preference dimensions
to guide product development and marketing strategy.

\textbf{Analysis Overview:}

\textbf{Preliminary Analysis:}

\begin{itemize}
\item
  KMO = 0.91 (Excellent)
\item
  Bartlett's test significant (p \textless{} 0.001)
\item
  No multicollinearity issues detected
\end{itemize}

\textbf{Factor Solution:}

\begin{itemize}
\item
  Method: Maximum Likelihood with Varimax rotation
\item
  Four factors extracted (eigenvalues \textgreater{} 1.0)
\item
  Total variance explained: 67.8\%
\end{itemize}

\textbf{Factor Interpretation:}

\textbf{Factor 1: Performance \& Technology (24.3\% variance)}

\begin{itemize}
\item
  Processor speed (0.84)
\item
  RAM capacity (0.79)
\item
  Graphics performance (0.76)
\item
  Operating system (0.71)
\item
  Storage capacity (0.68)
\end{itemize}

\textbf{Factor 2: Design \& Aesthetics (18.7\% variance)}

\begin{itemize}
\item
  Physical appearance (0.82)
\item
  Color options (0.78)
\item
  Build quality (0.75)
\item
  Weight and thickness (0.72)
\item
  Screen size (0.65)
\end{itemize}

\textbf{Factor 3: Value \& Practicality (13.9\% variance)}

\begin{itemize}
\item
  Price affordability (0.81)
\item
  Battery life (0.77)
\item
  Durability (0.74)
\item
  Warranty coverage (0.69)
\item
  Repair availability (0.63)
\end{itemize}

\textbf{Factor 4: Connectivity \& Features (10.9\% variance)}

\begin{itemize}
\item
  Camera quality (0.79)
\item
  Wireless capabilities (0.75)
\item
  Apps availability (0.71)
\item
  Internet speed (0.68)
\item
  Social media integration (0.64)
\end{itemize}

\textbf{Business Implications:}

\begin{itemize}
\item
  Four distinct consumer segments identified
\item
  Product positioning strategies developed for each factor
\item
  Marketing messages tailored to preference dimensions
\item
  Feature prioritization based on factor importance
\end{itemize}

\subsubsection{Case Study 3: Educational Research - Academic Motivation}

\textbf{Background:} Educational researchers administered a 36-item
Academic Motivation Scale to 650 university students to understand
different types of academic motivation.

\textbf{Theoretical Framework:} Self-Determination Theory predicts seven
types of motivation:

\begin{enumerate}
\item
  Intrinsic Motivation (3 types)
\item
  Extrinsic Motivation (3 types)
\end{enumerate}

\begin{enumerate}
\setcounter{enumi}{2}
\item
  Amotivation (1 type)
\end{enumerate}

\textbf{Analytical Challenge:} Test whether the data supports the
theoretical seven-factor structure versus alternative models.

\textbf{Model Comparison Approach:}

\textbf{Model 1: Seven-Factor Solution}

\begin{itemize}
\item
  Factors: IM-Knowledge, IM-Accomplishment, IM-Stimulation,
  EM-Identified, EM-Introjected, EM-External, Amotivation
\item
  Fit: χ² = 892.3, df = 573, CFI = 0.94, RMSEA = 0.058
\item
  All factor loadings \textgreater{} 0.50
\item
  Good interpretability
\end{itemize}

\textbf{Model 2: Three-Factor Solution}

\begin{itemize}
\item
  Factors: Intrinsic, Extrinsic, Amotivation
\item
  Fit: χ² = 1,247.8, df = 591, CFI = 0.89, RMSEA = 0.078
\item
  Some loss of discriminant validity
\end{itemize}

\textbf{Model 3: One-Factor Solution}

\begin{itemize}
\item
  Single ``Academic Motivation'' factor
\item
  Fit: χ² = 2,156.4, df = 594, CFI = 0.71, RMSEA = 0.122
\item
  Poor fit, not theoretically meaningful
\end{itemize}

\textbf{Results:} The seven-factor model provided the best fit to the
data, supporting the theoretical distinction between different types of
academic motivation.

\textbf{Factor Correlations (Oblique Rotation):}

\begin{verbatim}
                   IM-K  IM-A  IM-S  EM-I  EM-J  EM-E  AMOT
IM-Knowledge        1.00
IM-Accomplishment   0.73  1.00
IM-Stimulation      0.68  0.71  1.00
EM-Identified       0.45  0.52  0.41  1.00
EM-Introjected      0.21  0.28  0.18  0.65  1.00
EM-External         0.08  0.12  0.05  0.43  0.58  1.00
Amotivation        -0.35 -0.41 -0.38 -0.52 -0.28 -0.15  1.00
\end{verbatim}

The pattern shows expected relationships: intrinsic motivation types
correlate positively with each other and negatively with amotivation.

\subsubsection{Case Study 4: Health Research - Quality of Life
Assessment}

\textbf{Background:} Medical researchers developed a quality of life
questionnaire for cancer patients with 28 items across multiple life
domains.

\textbf{Methodological Considerations:}

\begin{itemize}
\item
  Small sample size (N = 185) relative to variables
\item
  Missing data due to patient fatigue
\item
  Ordinal rating scales (0-4)
\item
  Potential floor/ceiling effects
\end{itemize}

\textbf{Special Analytical Approaches:}

\textbf{Missing Data Handling:}

\begin{itemize}
\item
  Multiple imputation (m = 20 imputations)
\item
  Analysis conducted on each imputed dataset
\item
  Results pooled using Rubin's rules
\end{itemize}

\textbf{Ordinal Data Treatment:}

\begin{itemize}
\item
  Polychoric correlation matrix used
\item
  Weighted least squares estimation
\item
  Robust standard errors computed
\end{itemize}

\textbf{Factor Structure:} Four factors emerged representing different
quality of life domains:

\textbf{Factor 1: Physical Functioning}

\begin{itemize}
\item
  Physical symptoms and limitations
\item
  Energy and fatigue levels
\item
  Daily activity capacity
\end{itemize}

\textbf{Factor 2: Emotional Well-being}

\begin{itemize}
\item
  Anxiety and depression
\item
  Mood and emotional state
\item
  Coping with illness
\end{itemize}

\textbf{Factor 3: Social Functioning}

\begin{itemize}
\item
  Relationships with family/friends
\item
  Social activities participation
\item
  Support system quality
\end{itemize}

\textbf{Factor 4: Treatment Impact}

\begin{itemize}
\item
  Side effects management
\item
  Treatment satisfaction
\item
  Healthcare communication
\end{itemize}

\textbf{Validation Results:}

\begin{itemize}
\item
  Cross-validation with independent sample (N = 142)
\item
  Coefficient of congruence \textgreater{} 0.85 for all factors
\item
  Convergent validity with established QOL measures
\item
  Known-groups validity confirmed
\end{itemize}

\subsubsection{Common Challenges and Solutions}

\textbf{Challenge 1: Small Sample Sizes} \textbf{Solutions:}

\begin{itemize}
\item
  Use more conservative factor retention criteria
\item
  Focus on replication rather than exploration
\item
  Consider item parceling for complex models
\item
  Report confidence intervals for loadings
\end{itemize}

\textbf{Challenge 2: Mixed Data Types} \textbf{Solutions:}

\begin{itemize}
\item
  Use appropriate correlation matrices (polychoric, tetrachoric)
\item
  Consider categorical factor analysis methods
\item
  Transform variables when appropriate
\item
  Report sensitivity analyses
\end{itemize}

\textbf{Challenge 3: Cross-Loadings and Complex Structure}
\textbf{Solutions:}

\begin{itemize}
\item
  Try different rotation methods
\item
  Consider hierarchical factor models
\item
  Use confirmatory factor analysis
\item
  Examine item content for ambiguity
\end{itemize}

\textbf{Challenge 4: Non-Normal Data} \textbf{Solutions:}

\begin{itemize}
\item
  Use robust estimation methods
\item
  Consider non-parametric alternatives
\item
  Transform variables when appropriate
\item
  Report distributional assumptions
\end{itemize}

\textbf{Challenge 5: Interpretation Difficulties} \textbf{Solutions:}

\begin{itemize}
\item
  Examine multiple solutions (different numbers of factors)
\item
  Consider theoretical expectations
\item
  Use external validation criteria
\item
  Consult domain experts
\end{itemize}

---

\textbf{Study Questions for Section 2.4:}

\begin{enumerate}
\item
  In the personality case study, why was oblique rotation chosen over
  orthogonal rotation?
\item
  How would you handle a situation where parallel analysis suggests 3
  factors but theory predicts 5?
\item
  What are the advantages and disadvantages of using polychoric
  correlations for ordinal data?
\item
  How do you interpret factor correlations in an oblique solution?
\item
  What validation strategies are most appropriate for small sample
  factor analyses?
\item
  When might you choose a hierarchical factor model over a simple
  structure model?
\end{enumerate}

\subsection{Software Implementation and Advanced Topics}

This section covers practical software implementation of Factor Analysis
and advanced methodological considerations for complex research
scenarios.

\subsubsection{Software Packages and Implementation}

\textbf{Python Implementation:}

\textbf{Primary Libraries:}

\begin{itemize}
\item
  \texttt{factor\_analyzer}: Dedicated Factor Analysis library
\item
  \texttt{sklearn.decomposition}: PCA and Factor Analysis
\item
  \texttt{scipy.stats}: Statistical tests and correlations
\item
  \texttt{pandas}: Data manipulation
\item
  \texttt{numpy}: Numerical computations
\end{itemize}

\textbf{Basic Factor Analysis Workflow in Python:}

\begin{Shaded}
\begin{Highlighting}[]
\ImportTok{from}\NormalTok{ factor\_analyzer }\ImportTok{import}\NormalTok{ FactorAnalyzer}
\ImportTok{from}\NormalTok{ factor\_analyzer.factor\_analyzer }\ImportTok{import}\NormalTok{ calculate\_kmo, calculate\_bartlett\_sphericity}
\ImportTok{import}\NormalTok{ pandas }\ImportTok{as}\NormalTok{ pd}
\ImportTok{import}\NormalTok{ numpy }\ImportTok{as}\NormalTok{ np}

\CommentTok{\# Load and prepare data}
\NormalTok{data }\OperatorTok{=}\NormalTok{ pd.read\_csv(}\StringTok{\textquotesingle{}survey\_data.csv\textquotesingle{}}\NormalTok{)}
\NormalTok{X }\OperatorTok{=}\NormalTok{ data.select\_dtypes(include}\OperatorTok{=}\NormalTok{[np.number])}

\CommentTok{\# Test factorability}
\NormalTok{kmo\_all, kmo\_model }\OperatorTok{=}\NormalTok{ calculate\_kmo(X)}
\NormalTok{chi\_square\_value, p\_value }\OperatorTok{=}\NormalTok{ calculate\_bartlett\_sphericity(X)}

\BuiltInTok{print}\NormalTok{(}\SpecialStringTok{f\textquotesingle{}KMO: }\SpecialCharTok{\{}\NormalTok{kmo\_model}\SpecialCharTok{:.3f\}}\SpecialStringTok{\textquotesingle{}}\NormalTok{)}
\BuiltInTok{print}\NormalTok{(}\SpecialStringTok{f\textquotesingle{}Bartlett p{-}value: }\SpecialCharTok{\{}\NormalTok{p\_value}\SpecialCharTok{:.3f\}}\SpecialStringTok{\textquotesingle{}}\NormalTok{)}

\CommentTok{\# Determine number of factors}
\NormalTok{fa }\OperatorTok{=}\NormalTok{ FactorAnalyzer(rotation}\OperatorTok{=}\VariableTok{None}\NormalTok{, n\_factors}\OperatorTok{=}\NormalTok{X.shape[}\DecValTok{1}\NormalTok{])}
\NormalTok{fa.fit(X)}
\NormalTok{eigenvalues }\OperatorTok{=}\NormalTok{ fa.get\_eigenvalues()[}\DecValTok{0}\NormalTok{]}

\CommentTok{\# Perform Factor Analysis}
\NormalTok{fa\_final }\OperatorTok{=}\NormalTok{ FactorAnalyzer(n\_factors}\OperatorTok{=}\DecValTok{5}\NormalTok{, rotation}\OperatorTok{=}\StringTok{\textquotesingle{}promax\textquotesingle{}}\NormalTok{)}
\NormalTok{fa\_final.fit(X)}

\CommentTok{\# Extract results}
\NormalTok{loadings }\OperatorTok{=}\NormalTok{ fa\_final.loadings\_}
\NormalTok{communalities }\OperatorTok{=}\NormalTok{ fa\_final.get\_communalities()}
\NormalTok{factor\_variance }\OperatorTok{=}\NormalTok{ fa\_final.get\_factor\_variance()}
\end{Highlighting}
\end{Shaded}

\textbf{R Implementation:}

\textbf{Primary Packages:}

\begin{itemize}
\item
  \texttt{psych}: Comprehensive psychometric package
\item
  \texttt{GPArotation}: Advanced rotation methods
\item
  \texttt{corrplot}: Correlation visualization
\item
  \texttt{lavaan}: Confirmatory Factor Analysis
\end{itemize}

\textbf{Basic Factor Analysis in R:}

\begin{Shaded}
\begin{Highlighting}[]
\FunctionTok{library}\NormalTok{(psych)}
\FunctionTok{library}\NormalTok{(GPArotation)}

\CommentTok{\# Load data}
\NormalTok{data }\OtherTok{\textless{}{-}} \FunctionTok{read.csv}\NormalTok{(}\StringTok{"survey\_data.csv"}\NormalTok{)}

\CommentTok{\# Test factorability}
\FunctionTok{KMO}\NormalTok{(data)}
\FunctionTok{cortest.bartlett}\NormalTok{(}\FunctionTok{cor}\NormalTok{(data), }\AttributeTok{n =} \FunctionTok{nrow}\NormalTok{(data))}

\CommentTok{\# Determine number of factors}
\FunctionTok{fa.parallel}\NormalTok{(data, }\AttributeTok{fa =} \StringTok{"fa"}\NormalTok{)}
\FunctionTok{VSS}\NormalTok{(data, }\AttributeTok{rotate =} \StringTok{"varimax"}\NormalTok{)}

\CommentTok{\# Perform Factor Analysis}
\NormalTok{fa\_result }\OtherTok{\textless{}{-}} \FunctionTok{fa}\NormalTok{(data, }\AttributeTok{nfactors =} \DecValTok{5}\NormalTok{, }\AttributeTok{rotate =} \StringTok{"promax"}\NormalTok{, }\AttributeTok{fm =} \StringTok{"pa"}\NormalTok{)}

\CommentTok{\# Examine results}
\FunctionTok{print}\NormalTok{(fa\_result, }\AttributeTok{cut =} \FloatTok{0.3}\NormalTok{, }\AttributeTok{sort =} \ConstantTok{TRUE}\NormalTok{)}
\FunctionTok{fa.diagram}\NormalTok{(fa\_result)}
\end{Highlighting}
\end{Shaded}

\textbf{SPSS Implementation:}

\textbf{Menu Path: Analyze → Dimension Reduction → Factor}

\textbf{Key Options:}

\begin{itemize}
\item
  Extraction: Principal Axis Factoring or Maximum Likelihood
\item
  Rotation: Varimax (orthogonal) or Promax (oblique)
\item
  Scores: Save factor scores as variables
\item
  Options: KMO and Bartlett's test, communalities
\end{itemize}

\textbf{Critical SPSS Settings:}

\begin{itemize}
\item
  Extraction Criteria: Eigenvalues over 1, or fixed number
\item
  Maximum Iterations: Increase to 100 for convergence
\item
  Convergence: 0.001 for precise solutions
\item
  Display: Sorted by size, suppress small coefficients
\end{itemize}

\subsubsection{Advanced Factor Analysis Topics}

\subsubsection{Confirmatory Factor Analysis (CFA)}

\textbf{When to Use CFA:}

\begin{itemize}
\item
  Testing specific theoretical models
\item
  Validating measurement instruments
\item
  Cross-group comparisons
\item
  Longitudinal factor invariance
\end{itemize}

\textbf{CFA vs. Exploratory FA:}

\begin{verbatim}
Exploratory FA:
- Data-driven approach
- Discovers factor structure
- All variables can load on all factors
- Generates hypotheses

Confirmatory FA:
- Theory-driven approach
- Tests predetermined structure
- Specified loadings pattern
- Tests hypotheses
\end{verbatim}

\textbf{Model Specification in CFA:}

\begin{itemize}
\item
  Fixed parameters (typically loadings set to 0)
\item
  Free parameters (estimated from data)
\item
  Constraints (equality restrictions)
\item
  Identification requirements
\end{itemize}

\textbf{Fit Indices for CFA Models:}

\begin{itemize}
\item
  Chi-square test: Model vs. saturated model
\item
  CFI (Comparative Fit Index): ≥ 0.95 excellent
\item
  TLI (Tucker-Lewis Index): ≥ 0.95 excellent
\item
  RMSEA (Root Mean Square Error): \textless{} 0.06 excellent
\item
  SRMR (Standardized Root Mean Residual): \textless{} 0.08 good
\end{itemize}

\subsubsection{Hierarchical Factor Analysis}

\textbf{Concept:} Higher-order factors that explain correlations among
first-order factors.

\textbf{Model Structure:}

\begin{verbatim}
Second-Order Factor
    ├── First-Order Factor 1
    │   ├── Variable 1
    │   ├── Variable 2
    │   └── Variable 3
    ├── First-Order Factor 2
    │   ├── Variable 4
    │   ├── Variable 5
    │   └── Variable 6
    └── First-Order Factor 3
        ├── Variable 7
        ├── Variable 8
        └── Variable 9
\end{verbatim}

\textbf{Applications:}

\begin{itemize}
\item
  Intelligence research (g-factor)
\item
  Personality assessment (higher-order traits)
\item
  Quality of life measurement
\item
  Organizational behavior
\end{itemize}

\textbf{Advantages:}

\begin{itemize}
\item
  Parsimonious representation
\item
  Theoretical alignment
\item
  Explains factor correlations
\item
  Hierarchical interpretation
\end{itemize}

\subsubsection{Bifactor Models}

\textbf{Structure:} All variables load on a general factor plus specific
factors.

\textbf{Key Features:}

\begin{itemize}
\item
  General factor: Common to all variables
\item
  Specific factors: Orthogonal to general factor
\item
  No correlations between specific factors
\item
  Direct modeling of multidimensionality
\end{itemize}

\textbf{When to Use:}

\begin{itemize}
\item
  Unidimensionality assumption violated
\item
  Interest in general and specific components
\item
  Educational and psychological testing
\item
  Quality control applications
\end{itemize}

\textbf{Model Comparison:}

\begin{verbatim}
Traditional Model: F1 ↔ F2 ↔ F3 (correlated factors)
Hierarchical Model: G → F1, F2, F3 (higher-order)
Bifactor Model: G + S1, S2, S3 (general + specific)
\end{verbatim}

\subsubsection{Factor Analysis with Ordinal Data}

\textbf{Challenges with Ordinal Data:}

\begin{itemize}
\item
  Pearson correlations underestimate relationships
\item
  Distributional assumptions violated
\item
  Floor and ceiling effects
\item
  Limited response categories
\end{itemize}

\textbf{Polychoric Correlations:} Assume underlying continuous variables
for ordinal responses.

\textbf{Estimation Methods:}

\begin{itemize}
\item
  Weighted Least Squares (WLS)
\item
  Diagonally Weighted Least Squares (DWLS)
\item
  Robust Maximum Likelihood
\item
  Bayesian estimation
\end{itemize}

\textbf{Practical Considerations:}

\begin{itemize}
\item
  Minimum 5 response categories preferred
\item
  Check threshold parameters
\item
  Examine residuals carefully
\item
  Consider alternative models
\end{itemize}

\subsubsection{Missing Data in Factor Analysis}

\textbf{Missing Data Mechanisms:}

\begin{itemize}
\item
  MCAR (Missing Completely at Random)
\item
  MAR (Missing at Random)
\item
  MNAR (Missing Not at Random)
\end{itemize}

\textbf{Handling Strategies:}

\textbf{1. Listwise Deletion:}

\begin{itemize}
\item
  Simple but reduces sample size
\item
  Assumes MCAR mechanism
\item
  May introduce bias
\end{itemize}

\textbf{2. Pairwise Deletion:}

\begin{itemize}
\item
  Uses available data for each correlation
\item
  May produce non-positive definite matrices
\item
  Inconsistent sample sizes
\end{itemize}

\textbf{3. Multiple Imputation:}

\begin{itemize}
\item
  Creates multiple complete datasets
\item
  Accounts for uncertainty
\item
  Requires MAR assumption
\end{itemize}

\textbf{4. Full Information Maximum Likelihood (FIML):}

\begin{itemize}
\item
  Uses all available information
\item
  Efficient and unbiased under MAR
\item
  Available in SEM software
\end{itemize}

\textbf{Best Practices:}

\begin{itemize}
\item
  Examine missing data patterns
\item
  Test missing data assumptions
\item
  Use appropriate method for mechanism
\item
  Report sensitivity analyses
\end{itemize}

\subsubsection{Factor Score Estimation}

\textbf{Methods for Computing Factor Scores:}

\textbf{1. Regression Method:}
\[\hat{\mathbf{F}} = \mathbf{\Lambda}^{\top}\mathbf{R}^{- 1}\mathbf{X}\]

\begin{itemize}
\item
  Minimizes squared error
\item
  Correlated even for orthogonal factors
\item
  Most commonly used
\end{itemize}

\textbf{2. Bartlett Method:}
\[\hat{\mathbf{F}} = \left( \mathbf{\Lambda}^{\top}\mathbf{\Psi}^{- 1}\mathbf{\Lambda} \right)^{- 1}\mathbf{\Lambda}^{\top}\mathbf{\Psi}^{- 1}\mathbf{X}\]

\begin{itemize}
\item
  Unbiased estimates
\item
  Preserves factor correlations
\item
  Computationally intensive
\end{itemize}

\textbf{3. Anderson-Rubin Method:}

\begin{itemize}
\item
  Orthogonal scores
\item
  Unit variance
\item
  Zero correlations maintained
\end{itemize}

\textbf{Choosing Score Method:}

\begin{itemize}
\item
  Regression: General purpose, most interpretable
\item
  Bartlett: When factor correlations important
\item
  Anderson-Rubin: When orthogonality required
\end{itemize}

\textbf{Factor Score Applications:}

\begin{itemize}
\item
  Further statistical analyses
\item
  Group comparisons
\item
  Predictive modeling
\item
  Clustering analysis
\end{itemize}

\subsubsection{Advanced Rotation Methods}

\textbf{Gradient Projection Algorithms:}

\begin{itemize}
\item
  Faster convergence
\item
  Better local optima
\item
  Handles constraints efficiently
\end{itemize}

\textbf{Procrustes Rotation:}

\begin{itemize}
\item
  Rotates to target matrix
\item
  Cross-study comparisons
\item
  Theoretical alignment
\end{itemize}

\textbf{Independent Cluster Rotation:}

\begin{itemize}
\item
  Separate rotation of clusters
\item
  Complex factor structures
\item
  Exploratory-confirmatory hybrid
\end{itemize}

\textbf{Robust Rotation Methods:}

\begin{itemize}
\item
  Outlier-resistant
\item
  Non-normal data
\item
  Contaminated samples
\end{itemize}

\subsubsection{Model Selection and Comparison}

\textbf{Information Criteria:}

\begin{itemize}
\item
  AIC (Akaike Information Criterion)
\item
  BIC (Bayesian Information Criterion)
\item
  Sample-size adjusted BIC
\item
  CAIC (Consistent AIC)
\end{itemize}

\textbf{Cross-Validation Approaches:}

\begin{itemize}
\item
  Split-sample validation
\item
  K-fold cross-validation
\item
  Bootstrap validation
\item
  Permutation tests
\end{itemize}

\textbf{Nested Model Comparisons:}

\begin{itemize}
\item
  Likelihood ratio tests
\item
  Chi-square difference tests
\item
  Sequential model building
\item
  Parsimony considerations
\end{itemize}

\textbf{Non-Nested Comparisons:}

\begin{itemize}
\item
  Information criteria
\item
  Predictive accuracy
\item
  Theoretical coherence
\item
  Practical utility
\end{itemize}

\subsubsection{Contemporary Developments}

\textbf{Machine Learning Integration:}

\begin{itemize}
\item
  Factor Analysis as dimensionality reduction
\item
  Integration with clustering algorithms
\item
  Deep learning applications
\item
  Automated model selection
\end{itemize}

\textbf{Bayesian Factor Analysis:}

\begin{itemize}
\item
  Prior specification
\item
  Posterior distributions
\item
  Model uncertainty
\item
  Flexible modeling frameworks
\end{itemize}

\textbf{Network Psychometrics:}

\begin{itemize}
\item
  Alternative to latent variable models
\item
  Direct variable relationships
\item
  Partial correlation networks
\item
  Dynamic systems approach
\end{itemize}

\textbf{Big Data Considerations:}

\begin{itemize}
\item
  Scalable algorithms
\item
  Streaming factor analysis
\item
  Distributed computing
\item
  Approximate methods
\end{itemize}

---

\section{Conclusion and Summary}

This comprehensive study guide has covered the theoretical foundations,
practical implementation, and advanced applications of Principal
Component Analysis and Factor Analysis. These powerful multivariate
techniques provide researchers with tools to:

\begin{itemize}
\item
  Reduce data dimensionality while preserving essential information
\item
  Discover underlying latent structures in complex datasets
\item
  Validate theoretical models and measurement instruments
\item
  Extract meaningful patterns from large-scale surveys and assessments
\end{itemize}

\textbf{Key Takeaways:}

\textbf{Methodological Rigor:}

\begin{itemize}
\item
  Proper assessment of data suitability is crucial
\item
  Multiple criteria should guide factor retention decisions
\item
  Rotation method selection impacts interpretability
\item
  Validation is essential for reliable results
\end{itemize}

\textbf{Practical Application:}

\begin{itemize}
\item
  Software implementation requires understanding of underlying theory
\item
  Real-world data presents challenges requiring adaptive strategies
\item
  Domain expertise enhances statistical analysis
\item
  Communication of results must balance technical accuracy with
  accessibility
\end{itemize}

\textbf{Future Directions:}

\begin{itemize}
\item
  Integration with machine learning methods
\item
  Handling of big data and streaming applications
\item
  Development of robust methods for complex data
\item
  Advancement in Bayesian and network approaches
\end{itemize}

The field continues to evolve with new methodological developments and
computational advances, making these techniques increasingly powerful
and accessible for researchers across disciplines.

---

\textbf{Final Study Questions - Integration and Application:}

\begin{enumerate}
\item
  Compare and contrast the assumptions and applications of PCA versus
  Factor Analysis.
\item
  Design a complete factor analysis study for your research area,
  including power analysis and validation strategy.
\item
  How would you handle a dataset with mixed continuous and ordinal
  variables in factor analysis?
\item
  What are the trade-offs between exploratory and confirmatory
  approaches in factor analysis?
\item
  How do modern machine learning approaches complement traditional
  factor analysis methods?
\item
  What ethical considerations arise when using factor analysis in
  high-stakes assessment contexts?
\end{enumerate}

\textbf{Recommended Further Reading:}

\begin{itemize}
\item
  Fabrigar, L. R., \& Wegener, D. T. (2012). \textbf{Exploratory Factor
  Analysis}
\item
  Brown, T. A. (2015). \textbf{Confirmatory Factor Analysis for Applied
  Research}
\item
  Hayton, J. C., Allen, D. G., \& Scarpello, V. (2004). Factor retention
  decisions in exploratory factor analysis
\item
  Howard, M. C. (2016). A review of exploratory factor analysis
  decisions and overview of current practices
\end{itemize}

\end{document}

